%% MoCafe — Clumpy dust medium (dust-only).
%% Author: Kwang-Il Seon
%% Written 2026-05-30 for kiseon_memo.cls.
\documentclass[english,a4paper]{kiseon_memo}
\synctex=1
\usepackage{amsmath}
\usepackage{amssymb}
\usepackage{microtype}

\newcommand{\rd}{\mathrm{d}}
\newcommand{\ee}{\mathrm{e}}
\newcommand{\obs}{\mathrm{obs}}
\newcommand{\kc}{\kappa_{\mathrm c}}
\newcommand{\rc}{r_{\mathrm c}}
\newcommand{\fcov}{f_{\mathrm{cov}}}
\newcommand{\fvol}{f_{\mathrm{vol}}}
\newcommand{\Ncl}{N_{\mathrm{cl}}}

\begin{document}

\title{The Clumpy Dust Medium in MoCafe}
\author{Kwang-Il Seon}
\date{2026-05-30}
\maketitle

\begin{abstract}
MoCafe's clumpy-medium mode (\texttt{par\%grid\_type = 'clump'}) replaces the
Cartesian density grid with a population of $\Ncl$ non-overlapping spherical
dust clumps of radius $\rc$ placed inside a sphere of radius $R$; the
inter-clump space is vacuum.  The clumps are discrete spheres (not a two-phase
or continuous medium), implemented with the placement and ray-tracing machinery
ported from LaRT but stripped to grey dust radiative transfer:
because the dust opacity is frequency-independent, no Doppler width, Voigt
profile, frequency tracking, or velocity field is needed, and the clump
subsystem reduces to pure geometry plus a per-clump scalar opacity.  The
scattering and peeling-off physics are shared verbatim with the Cartesian
dust path; only the optical-depth integration differs.  This memo records the
model, the clump-count and opacity calibration, the Random-Sequential-Addition
placement, the CSR/DDA accelerated ray-tracer, and the validation against a
homogeneous sphere.
\end{abstract}

\section{Model}
A point or extended source illuminates $\Ncl$ spherical dust clumps, each of
radius $\rc$ (possibly varying with clump-center radius via a profile) and grey
dust opacity $\kc$ (extinction per unit path length).  Clump centers lie in the
spherical shell $r_{\min}\le r\le R$, where $R=$ \texttt{par\%rmax} and
$r_{\min}=\max(0,\texttt{par\%rmin})$; with $r_{\min}=0$ the shell is a full
ball.  Outside the clumps the medium is vacuum, so a photon alternates between
\emph{vacuum flight} (find the next clump it enters) and \emph{clump traversal}
(accumulate $\kc\,\ell$ along the chord of length $\ell$).  No opacity, and
hence no interaction, occurs between clumps.

\section{Clump count}
For uniform clumps the realized covering and volume filling factors of $\Ncl$
clumps in the shell are
\begin{equation}
  \fvol = \frac{\Ncl\,\rc^{3}}{R^{3}-r_{\min}^{3}},
  \qquad
  \fcov = \frac{3}{4}\,
          \frac{\Ncl\,\rc^{2}}{R^{2}+R\,r_{\min}+r_{\min}^{2}} ,
  \label{eq:fvolfcov}
\end{equation}
so that $\fvol = \tfrac{4}{3}\,\fcov\,\rc/R$ for $r_{\min}=0$.  The user
specifies exactly one of \texttt{clump\_N\_clumps}, \texttt{clump\_f\_vol}, or
\texttt{clump\_f\_cov}; the other two follow from \eqref{eq:fvolfcov}.  When a
radial \emph{number} profile $n_{\mathrm{cl}}(r)\propto$ \texttt{shape}$(r)$ is
active, $\Ncl$, $\fvol$, and $\fcov$ are obtained by quadrature of
$n_{\mathrm{cl}}(r)$ over the shell (the line-of-sight definition of $\fcov$).
Reliable random placement requires $\fvol\lesssim0.35$ (three-dimensional
jamming $\approx0.38$); MoCafe aborts with guidance if the placement stalls.

\section{Per-clump opacity}
The grey dust opacity $\kc$ of a clump is set, in priority order, by one of
\begin{equation}
  \kc = \frac{\tau_{0}}{\rc}
  \quad(\texttt{clump\_tau0}),
  \qquad
  \kc = n\,C_{\mathrm{ext}}\,\ell_{\mathrm{cm}}
  \quad(\texttt{clump\_ndust} \text{ or } \texttt{clump\_nH}),
  \label{eq:kappa}
\end{equation}
where $\tau_{0}$ is the dimensionless dust optical depth from a clump's center
to its surface, $C_{\mathrm{ext}}=$ \texttt{par\%cext\_dust} is the dust
extinction per H, and $\ell_{\mathrm{cm}}=$ \texttt{par\%distance2cm} converts
code length to cm (the same $n\,C_{\mathrm{ext}}\,\ell_{\mathrm{cm}}$ used by the
Cartesian grid).  Alternatively the population is calibrated to a system-level
radial target $\tau_{\max}$ or $\tau_{\mathrm{homo}}$ by back-solving
\begin{equation}
  \kc = \frac{\tau_{\mathrm{target}}}{\mathrm{GF}},\qquad
  \mathrm{GF} = \frac{\Ncl\,\rc^{3}}{R^{2}+R\,r_{\min}+r_{\min}^{2}}
  \quad(\text{uniform}),
\end{equation}
with the profile-mode geometric factor obtained by quadrature.

\paragraph{System scalars.}
From the placed population MoCafe reports
\begin{align}
  \tau_{\mathrm{homo}} &= \sum_{i}
     \frac{\kc^{(i)}\,\rc^{(i)\,3}}{R^{2}+R\,r_{\min}+r_{\min}^{2}},
  \label{eq:tauhomo}\\
  \tau_{\max} &= \sum_{i}
     \frac{\kc^{(i)}\,\rc^{(i)\,3}}{3\,\max\!\big(d_{i}^{2},\,\rc^{(i)\,2}\big)},
  \label{eq:taumax}
\end{align}
where $d_{i}=|\mathbf{r}_{i}|$ is the clump-center distance from the origin.
Equation~\eqref{eq:tauhomo} is the opacity of the clumps smeared uniformly over
the shell and traversed radially; \eqref{eq:taumax} is the small-angle expected
radial optical depth (the $\max$ caps the singularity for a clump straddling
the origin).  In the uniform-isotropic limit both reduce to
$\Ncl\,\rc^{3}\kc/(R^{2}+R\,r_{\min}+r_{\min}^{2}) = \tfrac{4}{3}\,\fcov\,\tau_0$
(the uniform discrete-clump closed form).

\section{Placement: Random Sequential Addition}
Clump centers are drawn one at a time and accepted only if the new clump does
not overlap any previously placed one (\textsc{rsa}).  A linked-list hash grid
with cells of side $\ge 2\,\rc^{\max}$ makes the overlap test $O(1)$: a trial
center is checked only against clumps registered in its own cell and the 26
neighbors (the 27-cell rule).  Uniform placement uses box rejection inside the
shell; profile placement draws the center radius from the tabulated inverse CDF
of $n_{\mathrm{cl}}(r)\,r^{2}$ and isotropic (or biconical) angles.  Placement
runs once on the global root rank and is broadcast to one rank per node, then
shared within the node through MPI-3 shared memory, so every node holds an
identical clump layout.

\section{Accelerated ray-tracing}
\paragraph{CSR acceleration grid.}
A uniform $C^{3}$ grid over the bounding box indexes the clumps in
compressed-sparse-row form: a two-pass build fills \texttt{cg\_start} (per-cell
offsets) and \texttt{cg\_list} (clump indices), registering each clump in every
cell its bounding box overlaps.  The grid and the clump arrays live in MPI-3
shared memory (one copy per node).

\paragraph{Ray--sphere intersection.}
For a ray $P(t)=\mathbf{o}+t\,\hat{\mathbf k}$ and a clump centered at
$\mathbf c$ with radius $\rc$, writing $\mathbf{b}=\mathbf{o}-\mathbf{c}$,
\begin{equation}
  t_{\pm} = -(\mathbf{b}\!\cdot\!\hat{\mathbf k})
            \pm\sqrt{(\mathbf{b}\!\cdot\!\hat{\mathbf k})^{2}
                     -|\mathbf{b}|^{2}+\rc^{2}} ,
\end{equation}
with a real intersection when the discriminant is non-negative and $t_{+}>0$.

\paragraph{DDA traversal.}
The next clump a photon in vacuum will enter is found by stepping the ray
through the CSR grid with the Amanatides \& Woo (1987) digital differential
analyzer: the per-axis distances $t_{x},t_{y},t_{z}$ to the next cell face are
advanced by $\Delta t_{x}=\Delta x/|k_{x}|$ (and similarly for $y,z$); at each
cell every registered clump is tested with the ray--sphere formula and the
nearest positive entry is kept.  Traversal stops once the cell-entry distance
exceeds the best hit or the sphere boundary.

\paragraph{Transport.}
\texttt{raytrace\_to\_tau\_clump} advances the photon until a sampled optical
depth $\tau$ is accumulated: inside a clump it scatters at
$\rd s=\tau_{\mathrm{rem}}/\kc$ if $\tau_{\mathrm{rem}}\le\kc\,\ell$, otherwise
it subtracts $\kc\,\ell$, steps to the clump exit, and resumes in vacuum;
in vacuum it advances to the next clump (or escapes through the sphere).
\texttt{raytrace\_to\_edge\_clump} sums $\kc\,\ell$ over every clump the
line of sight crosses to the sphere boundary, and is used by the peeling-off
estimator.

\section{Reuse of the dust physics}
Clump mode swaps only the two ray-tracing procedure pointers
(\texttt{raytrace\_to\_tau}, \texttt{raytrace\_to\_edge}); the forced first
scattering, the Henyey--Greenstein (or Mueller-matrix) scattering, and the
peeling-off to each observer are \emph{identical} to the Cartesian dust path.
This is possible because the peel routines integrate the optical depth to the
observer through the \texttt{raytrace\_to\_edge} pointer rather than calling the
Cartesian routine directly; the default binding keeps the Cartesian results
bit-identical.  The forced-first-scattering weight $1-\ee^{-\tau_{0}}$ uses the
clump optical depth $\tau_{0}$ to the sphere edge along the emission direction,
which is zero (the photon escapes unscattered) for sight lines that miss every
clump.  The birth clump is identified at emission so a source inside a clump is
handled correctly.

\section{File I/O and the standalone generator}
With \texttt{par\%save\_clump\_info} the population is written to a binary
table (FITS or HDF5): columns \texttt{X, Y, Z} (centers) plus \texttt{R\_CLUMP}
and \texttt{RHOKAP} when non-uniform, and header keywords
\texttt{N\_CLUMPS, SPHERE\_R, RMIN, CL\_RAD, F\_VOL, F\_COV, TAU0, RHOKAP,
RMAX, DISTUNIT, DIST\_CM}.  A pre-built population is read back through
\texttt{par\%clump\_input\_file} (a \texttt{DENSITY}/\texttt{DENS} column is
converted to opacity with $C_{\mathrm{ext}}\,\ell_{\mathrm{cm}}$), optionally
rescaled so its realized $\tau_{\max}$ or $\tau_{\mathrm{homo}}$ matches a
target.  The Python helper \texttt{python/make\_clumps.py} writes such files
(uniform random populations or a single centered clump).

\section{Validation}
\paragraph{Homogeneous limit.}
A single clump centered at the origin with $\rc=R$ and $\kc=\tau_{0}/R$ fills
the sphere with uniform opacity, so it must reproduce the homogeneous dust
sphere of radial optical depth $\tau_{0}$.  At $\tau_{0}=1$, albedo $0.5$,
isotropic scattering, $2\times10^{6}$ photons, the integrated images agree with
the Cartesian reference to Monte-Carlo precision:
\begin{center}
\begin{tabular}{lccc}
\hline
section    & homogeneous & giant clump & rel.\ diff.\\
\hline
Scattered  & $6.38690\times10^{1}$ & $6.38682\times10^{1}$ & $-0.001\%$\\
Direct     & $1.21791\times10^{2}$ & $1.21791\times10^{2}$ & $+0.000\%$\\
Direct0    & $3.31062\times10^{2}$ & $3.31062\times10^{2}$ & $+0.000\%$\\
\hline
\end{tabular}
\end{center}
The direct images are deterministic (one source point) and match exactly; the
scattered totals match to $10^{-5}$.  The mean number of scatterings agrees to
$0.06\%$.  More generally the clumpy result approaches the smooth sphere as
$\fcov$ grows at fixed $\tau_{\mathrm{homo}}$.

\section{Scope}
Gas velocity and temperature are out of scope at this stage, so the LaRT
Doppler width, Voigt profile, frequency tracking, per-clump velocity, and the
Lyman-$\alpha$ resonance machinery are not ported.  The ray-tracer assumes
non-overlapping clumps (RSA enforces this); the overlap-aware path is recorded
for a future revision.  See \texttt{AMR\_CLUMPS\_PLAN.md} Part~B.

\begin{thebibliography}{9}
\bibitem{AW87} Amanatides, J., \& Woo, A. 1987, A Fast Voxel Traversal
  Algorithm for Ray Tracing, in Eurographics '87, 3
\end{thebibliography}

\end{document}
